\documentclass[10pt,a4paper]{amsart}
\usepackage{amsmath,amssymb,amsthm,graphicx}
% Retained minimal preamble for the Stacks Project chapter flat.tex (More on Flatness).
% Verbatim excerpt of the chapter preamble: Xy-pic setup (native preamble lines 14-18) and the native
% theorem declarations with the mathematical macros (native preamble lines 154-217). Copied mechanically;
% nothing added, removed or redefined. The wrapper class and the amsmath/amssymb/amsthm/graphicx base
% packages are supplied by the assembler, so the class/externaldocument part of the native preamble is omitted.
% For commutative diagrams we use Xy-pic
\usepackage[all]{xy}

% We use 2cell for 2-commutative diagrams.
\xyoption{2cell}
% Theorem environments.
%
\theoremstyle{plain}
\newtheorem{theorem}[subsection]{Theorem}
\newtheorem{proposition}[subsection]{Proposition}
\newtheorem{lemma}[subsection]{Lemma}

\theoremstyle{definition}
\newtheorem{definition}[subsection]{Definition}
\newtheorem{example}[subsection]{Example}
\newtheorem{exercise}[subsection]{Exercise}
\newtheorem{situation}[subsection]{Situation}

\theoremstyle{remark}
\newtheorem{remark}[subsection]{Remark}
\newtheorem{remarks}[subsection]{Remarks}

\numberwithin{equation}{subsection}

% Macros
%
\def\lim{\mathop{\mathrm{lim}}\nolimits}
\def\colim{\mathop{\mathrm{colim}}\nolimits}
\def\Spec{\mathop{\mathrm{Spec}}}
\def\Hom{\mathop{\mathrm{Hom}}\nolimits}
\def\Ext{\mathop{\mathrm{Ext}}\nolimits}
\def\SheafHom{\mathop{\mathcal{H}\!\mathit{om}}\nolimits}
\def\SheafExt{\mathop{\mathcal{E}\!\mathit{xt}}\nolimits}
\def\Sch{\mathit{Sch}}
\def\Mor{\mathop{\mathrm{Mor}}\nolimits}
\def\Ob{\mathop{\mathrm{Ob}}\nolimits}
\def\Sh{\mathop{\mathit{Sh}}\nolimits}
\def\NL{\mathop{N\!L}\nolimits}
\def\CH{\mathop{\mathrm{CH}}\nolimits}
\def\proetale{{pro\text{-}\acute{e}tale}}
\def\etale{{\acute{e}tale}}
\def\QCoh{\mathit{QCoh}}
\def\Ker{\mathop{\mathrm{Ker}}}
\def\Im{\mathop{\mathrm{Im}}}
\def\Coker{\mathop{\mathrm{Coker}}}
\def\Coim{\mathop{\mathrm{Coim}}}

% Boxtimes
%
\DeclareMathSymbol{\boxtimes}{\mathbin}{AMSa}{"02}

%
% Macros for moduli stacks/spaces
%
\def\QCohstack{\mathcal{QC}\!\mathit{oh}}
\def\Cohstack{\mathcal{C}\!\mathit{oh}}
\def\Spacesstack{\mathcal{S}\!\mathit{paces}}
\def\Quotfunctor{\mathrm{Quot}}
\def\Hilbfunctor{\mathrm{Hilb}}
\def\Curvesstack{\mathcal{C}\!\mathit{urves}}
\def\Polarizedstack{\mathcal{P}\!\mathit{olarized}}
\def\Complexesstack{\mathcal{C}\!\mathit{omplexes}}
% \Pic is the operator that assigns to X its picard group, usage \Pic(X)
% \Picardstack_{X/B} denotes the Picard stack of X over B
% \Picardfunctor_{X/B} denotes the Picard functor of X over B
\def\Pic{\mathop{\mathrm{Pic}}\nolimits}
\def\Picardstack{\mathcal{P}\!\mathit{ic}}
\def\Picardfunctor{\mathrm{Pic}}
\def\Deformationcategory{\mathcal{D}\!\mathit{ef}}

% --- appended 2026-10-12 (Chief): the project own definitions of the reference/slogan/history
% display environments, copied mechanically from the shared Stacks preamble lines 8-12
% (\usepackage{verbatim} plus the three \newenvironment lines, which the chapter preamble keeps
% verbatim from the project preamble). Nothing added, removed or redefined: these environments
% carry no mathematical content and the project itself typesets them as comments. Without these
% three lines a counted statement whose own text uses them fails with Environment slogan undefined.
\usepackage{verbatim}
\newenvironment{reference}{\comment}{\endcomment}
\newenvironment{slogan}{\comment}{\endcomment}
\newenvironment{history}{\comment}{\endcomment}
\usepackage{cleveref}
\crefname{originalequation}{original equation}{original equations}
\Crefname{originalequation}{Original equation}{Original equations}
\makeatletter
\newlabel{equation-star}{{0AT6}{source}{}{sourcepointer-1}{}}
\newlabel{equation-star@cref}{{[originalequation][0][]0AT6}{[1][0][]source}{}{}{}}
\makeatother
\crefname{originallemma}{original lemma}{original lemmas}
\Crefname{originallemma}{Original lemma}{Original lemmas}
\makeatletter
\newlabel{more-algebra-lemma-content-exists-flat-Mittag-Leffler}{{0ASD}{source}{}{sourcepointer-2}{}}
\newlabel{more-algebra-lemma-content-exists-flat-Mittag-Leffler@cref}{{[originallemma][0][]0ASD}{[1][0][]source}{}{}{}}
\makeatother
\crefname{originalsection}{original section}{original sections}
\Crefname{originalsection}{Original section}{Original sections}
\makeatletter
\newlabel{algebra-example-ML}{{059R}{source}{}{sourcepointer-3}{}}
\newlabel{algebra-example-ML@cref}{{[originalsection][0][]059R}{[1][0][]source}{}{}{}}
\makeatother
\crefname{originallemma}{original lemma}{original lemmas}
\Crefname{originallemma}{Original lemma}{Original lemmas}
\makeatletter
\newlabel{lemma-universally-injective-local}{{05FQ}{source}{}{sourcepointer-4}{}}
\newlabel{lemma-universally-injective-local@cref}{{[originallemma][0][]05FQ}{[1][0][]source}{}{}{}}
\makeatother
\title{Existence of a pure spreadout gives content ideals and kills a counterexample to universal injectivity}
\author{The Stacks Project Authors}
\date{stacks-project:flat.tex:local-snapshot-20261011; GNU FDL 1.2 or later}
\begin{document}
\maketitle
\noindent\textit{Editorial note.} All delivered mathematical text is the author's own, copied from the pinned Stacks Project chapter file flat.tex; no statement or proof step was rewritten, reordered or removed.\par\medskip
\noindent\textit{Editorial note.} The preamble is the chapter's own retained minimal preamble (sources/stacks-flat/retained-preamble.tex), a verbatim excerpt of the native Xy-pic setup and of the native theorem declarations and mathematical macros, followed by the project own definitions of the reference, slogan and history display environments (verbatim package plus three newenvironment lines) that the retained file carries at its end and that the project itself typesets as comments. No macro was invented or redefined, and the wrapper class and the amsmath/amssymb/amsthm/graphicx base packages are supplied by the assembler.\par\medskip
\noindent\textit{Editorial note.} The counted statement and proof contain no cross reference, no citation and no included file, so the delivered body is self contained.\par\medskip
\noindent\textit{Editorial note.} Canonical location of the counted proposition: Stacks tag 0AT7 (sources/stacks-flat/pinned-original-source.tex line 7229).\par\medskip
\noindent\textit{Allowed standard prerequisites (editorial).}\begin{itemize}
\item Pure spreadouts and the defining equation. Same chapter, the definition of a pure spreadout preceding this lemma and equation labelled equation-star, invoked in the statement. Supplies U, N' and the A-universal injectivity of N' -> N.
\item Existence of content ideals for projectively presented modules. More on Algebra, Lemma labelled more-algebra-lemma-content-exists-flat-Mittag-Leffler, cited at proof lines 7246-7248, with Algebra, Example labelled algebra-example-ML. Gives content ideals of elements of N' and hence of elements of N.
\item Universal injectivity modulo the maximal ideal implies universal injectivity. Same chapter, Lemma labelled lemma-universally-injective-local, cited at proof lines 7261-7263. Makes each composed map universally injective and then the colimit.
\end{itemize}\medskip
\ifdefined\hypertarget\hypertarget{sourcepointer-1}{}\fi
\noindent\textit{Source reference.} Original equation 0AT6: sources/stacks-flat/pinned-original-source.tex, source lines 7208--7208. Same-chapter pointer into the author environment equation at source line 7208; invoked at proof line None of the counted proof. It is retained only as a navigation pointer; every load-bearing step is inside the delivered span.\par\medskip
\ifdefined\hypertarget\hypertarget{sourcepointer-2}{}\fi
\noindent\textit{Source reference.} Original lemma 0ASD: sources/stacks-chow/tags.txt, source lines 13576--13576. Cross-chapter standard prerequisite located in the project tag list staged in the local snapshot; invoked at proof line 7245 of the counted proof. It is retained only as a navigation pointer; every load-bearing step is inside the delivered span.\par\medskip
\ifdefined\hypertarget\hypertarget{sourcepointer-3}{}\fi
\noindent\textit{Source reference.} Original section 059R: sources/stacks-chow/tags.txt, source lines 6824--6824. Cross-chapter standard prerequisite located in the project tag list staged in the local snapshot; invoked at proof line 7246 of the counted proof. It is retained only as a navigation pointer; every load-bearing step is inside the delivered span.\par\medskip
\ifdefined\hypertarget\hypertarget{sourcepointer-4}{}\fi
\noindent\textit{Source reference.} Original lemma 05FQ: sources/stacks-flat/pinned-original-source.tex, source lines 1946--1946. Same-chapter pointer into the author environment lemma at source line 1946; invoked at proof line 7264 of the counted proof. It is retained only as a navigation pointer; every load-bearing step is inside the delivered span.\par\medskip
\begin{lemma}
\label{lemma-properties-pure-spreadout}
In (\ref{equation-star}) if there exists a pure spreadout, then
\begin{enumerate}
\item elements of $N$ have content ideals in $A$, and
\item if $u : N \to M$ is a morphism to a flat $A$-module $M$
such that $N/\mathfrak m N \to M/\mathfrak m M$ is injective
for all maximal ideals $\mathfrak m$ of $A$, then $u$ is
$A$-universally injective.
\end{enumerate}
\end{lemma}

\begin{proof}
Choose $U$, $N'$ as in the definition of a pure spreadout.
Any element $x' \in N'$ has a content ideal in $A$ because
$N'$ is $A$-projective (this can easily be seen directly, but
it also follows from More on Algebra, Lemma
\ref{more-algebra-lemma-content-exists-flat-Mittag-Leffler} and
Algebra, Example \ref{algebra-example-ML}).
Since $N' \to N$ is $A$-universally injective, we see that
the image $x \in N$ of any $x' \in N'$ has a content ideal in $A$
(it is the same as the content ideal of $x'$).
For a general $x \in N$ we choose $s \in S$ such that
$s x$ is in the image of $N' \to N$ and we use that $x$ and $sx$
have the same content ideal.

\medskip\noindent
Let $u : N \to M$ be as in (2). To show that $u$ is $A$-universally
injective, we may replace $A$ by a localization at a maximal ideal
(small detail omitted). Assume $A$ is local with maximal ideal $\mathfrak m$.
Pick $s \in S$ and consider the composition
$$
N' \to N \xrightarrow{1/s} N \xrightarrow{u} M
$$
Each of these maps is injective modulo $\mathfrak m$, hence the composition
is $A$-universally injective by
Lemma \ref{lemma-universally-injective-local}.
Since $N = \colim_{s \in S} (1/s)N'$ we conclude that $u$
is $A$-inversally injective as a colimit of universally injective maps.
\end{proof}

\end{document}
